% Options for packages loaded elsewhere
\PassOptionsToPackage{unicode}{hyperref}
\PassOptionsToPackage{hyphens}{url}
\documentclass[
]{ltjsarticle}
\usepackage{xcolor}
\usepackage{amsmath,amssymb}
\setcounter{secnumdepth}{-\maxdimen} % remove section numbering
\usepackage{iftex}
\ifPDFTeX
  \usepackage[T1]{fontenc}
  \usepackage[utf8]{inputenc}
  \usepackage{textcomp} % provide euro and other symbols
\else % if luatex or xetex
  \usepackage{unicode-math} % this also loads fontspec
  \defaultfontfeatures{Scale=MatchLowercase}
  \defaultfontfeatures[\rmfamily]{Ligatures=TeX,Scale=1}
\fi
\usepackage{lmodern}
\ifPDFTeX\else
  % xetex/luatex font selection
\fi
% Use upquote if available, for straight quotes in verbatim environments
\IfFileExists{upquote.sty}{\usepackage{upquote}}{}
\IfFileExists{microtype.sty}{% use microtype if available
  \usepackage[]{microtype}
  \UseMicrotypeSet[protrusion]{basicmath} % disable protrusion for tt fonts
}{}
\makeatletter
\@ifundefined{KOMAClassName}{% if non-KOMA class
  \IfFileExists{parskip.sty}{%
    \usepackage{parskip}
  }{% else
    \setlength{\parindent}{0pt}
    \setlength{\parskip}{6pt plus 2pt minus 1pt}}
}{% if KOMA class
  \KOMAoptions{parskip=half}}
\makeatother
\setlength{\emergencystretch}{3em} % prevent overfull lines
\providecommand{\tightlist}{%
  \setlength{\itemsep}{0pt}\setlength{\parskip}{0pt}}
\usepackage{newunicodechar}
\newunicodechar{→}{\ensuremath{\rightarrow}}
\newunicodechar{⟨}{\ensuremath{\langle}}
\newunicodechar{⟩}{\ensuremath{\rangle}}
\newunicodechar{⇒}{\ensuremath{\Rightarrow}}
\newunicodechar{⟹}{\ensuremath{\Longrightarrow}}
\newunicodechar{×}{\ensuremath{\times}}
\newunicodechar{≈}{\ensuremath{\approx}}
\newunicodechar{∼}{\ensuremath{\sim}}
\newunicodechar{≤}{\ensuremath{\le}}
\newunicodechar{≥}{\ensuremath{\ge}}
\newunicodechar{≠}{\ensuremath{\ne}}
\newunicodechar{∫}{\ensuremath{\int}}
\newunicodechar{−}{\ensuremath{-}}
\newunicodechar{✗}{\ensuremath{\times}}
\newunicodechar{✓}{\ensuremath{\checkmark}}
\usepackage{bookmark}
\IfFileExists{xurl.sty}{\usepackage{xurl}}{} % add URL line breaks if available
\urlstyle{same}
\hypersetup{
  hidelinks,
  pdfcreator={LaTeX via pandoc}}

\author{}
\date{}

\begin{document}

\section{観測者--系ビート予想の関連先行研究調査 ―
先行プログラムとの関係および検証可能領域の所在（補遺）}\label{ux89b3ux6e2cux8005ux7cfbux30d3ux30fcux30c8ux4e88ux60f3ux306eux95a2ux9023ux5148ux884cux7814ux7a76ux8abfux67fb-ux5148ux884cux30d7ux30edux30b0ux30e9ux30e0ux3068ux306eux95a2ux4fc2ux304aux3088ux3073ux691cux8a3cux53efux80fdux9818ux57dfux306eux6240ux5728ux88dcux907a}

\textbf{副題}:
前論文公開後に行った事後的な文献調査の記録。関連の見込まれる先行研究を必要部分のみ引用し、前論文との接続は今後の課題として残す。

\textbf{木原 範昭} （本稿は前論文 {[}1{]}
の\textbf{補遺}である。確立物理の主張ではなく、新たな数式の導出も物理的主張も含まない。前論文
{[}1{]}
の公開後に、その機構と関連する可能性のある先行研究を事後的に調査した結果を記録するものである。ここで挙げる文献が前論文の予想と論理的に接続するか否かは確定しておらず、\textbf{今後の課題}である。本稿はそれらが前論文の予想を裏づける証拠であるとは主張せず、関連性が見込まれる範囲で必要部分のみを引用し、どの点が近く・どの点が異なるかを整理するにとどめる。）

バージョン v0.1（2026-06-28） DOI (Version): 10.5281/zenodo.20985750
（本版 v0.1） DOI (Concept):
10.5281/zenodo.20985749（外部参照用・最新版へ自動転送） Zenodo:
https://zenodo.org/records/20985750

\begin{center}\rule{0.5\linewidth}{0.5pt}\end{center}

\subsection{要旨}\label{ux8981ux65e8}

前論文 {[}1{]}（以下「ビート予想」）は、局在核論文 {[}2{]}
が公準として残した「ランダム性の起源」に対する機構の候補として、観測器--系の相対位相（ビート／ヘテロダイン）が無理振動数比のとき
Weyl
等分布することを提示し、再生核性とあわせて位置分解強度プロファイルの形が
\(|\psi_{\rm base}|^2\)
に一致すること（核）と、強度から単発クリックへの移行（橋）を区別した。前論文は
v0.4 において、有限 \(N\) の \(O(1/N)\)
補正を\textbf{測定可能な逸脱とは主張せず}、ビート機構を不確定性の一面と位置づけ、より根源的な不確定性が完全コヒーレント（厳密な整数比）な観測の実現そのものを妨げる――そのような機構が存在する可能性があるという仮説を置き（その存在は主張しない）、もしそれがあれば標準理論物理学のボルン則が\textbf{より強く成立する}、という立場をとる
{[}1{]}。

本補遺は、ビート予想の公開後に行った関連文献調査の記録である。調査の結論は二点である。第一に、\textbf{ビート予想の中核機構（系--観測器の振動数ビート＋無理比の
Weyl
等分布によるボルン統計の創発、およびコヒーレント／整合比領域でのボルン則からの逸脱）を、同じ形で提案・検証・観測した先行研究は、調査した範囲では確認されなかった}。第二に、それでもなお、\textbf{「ボルン則を決定論的な基層と閾値検出から創発／近似として導く」という問題意識を共有しうる先行プログラムが複数存在し}、ビート予想はそれらと関連する可能性がある。本補遺はこれらを
(1) 同じ問題意識をもちうる先行プログラム、(2)
別方向の研究、に分けて整理し、あわせて、ビート予想を検証しうる精密分光（原子時計）の領域が、判別に必要な測定（固定相対位相での単発射影統計を整合比の関数として見る測定）を、確認した範囲では行っていないことを記録する。いずれの引用についても、機構の同一性は主張しない。

\begin{center}\rule{0.5\linewidth}{0.5pt}\end{center}

\subsection{1.
はじめに：補遺の動機と位置づけ}\label{ux306fux3058ux3081ux306bux88dcux907aux306eux52d5ux6a5fux3068ux4f4dux7f6eux3065ux3051}

ビート予想 {[}1{]}
は予想（予言）論文であり、確立物理の主張ではない。その公開後、本予想が既存研究の系譜に対してどこに位置するかを正確に把握するため、関連の見込まれる先行研究を事後的に調査した。

この作業の目的は二つである。第一に、誠実さの担保である。新規の枠組みを既存文献に接地させずに提示することは、主張の射程を見誤らせる。第二に、本予想の独自性（あるとすれば何か）を、既知の系譜との差分として正確に切り出すことである。

本補遺は、調査の結果を記録するにとどまり、次を\textbf{主張しない}。

\begin{itemize}
\tightlist
\item
  引用する先行研究がビート予想を裏づける、とは主張しない。
\item
  引用する先行研究とビート予想が同一の機構である、とは主張しない。
\item
  いずれかの実験がビート予想を検証した、とは主張しない。
\end{itemize}

引用は、関連性が見込まれる必要部分に限る。前論文 {[}1{]}
が既に引用している文献（{[}6{]}{[}8{]}
など）については、本補遺で内容を再確認したうえで、本予想との関連可能性を改めて整理する。

\begin{center}\rule{0.5\linewidth}{0.5pt}\end{center}

\subsection{2.
調査の方法と射程}\label{ux8abfux67fbux306eux65b9ux6cd5ux3068ux5c04ux7a0b}

調査は、(i) ボルン則を決定論的基層・古典場・閾値検出から導く
foundational な試み、(ii)
コヒーレントな局部発振器を用いるボルン則検定、(iii)
原子時計におけるロック点近傍の量子射影雑音（QPN）の振る舞いと、エンタングルメントによらない
sub-QPN 異常の有無、(iv)
任意関数の方法のボルン則への適用、の四領域を中心に行った。各候補文献は一次資料（出版版または
arXiv 原文）にあたって内容を確認した。

本調査は網羅的ではない。確認したのは有限個の文献であり、「該当する先行研究が存在しない」という結論は、調査した範囲についての記述であって、文献全体にわたる不在の証明ではない。

\begin{center}\rule{0.5\linewidth}{0.5pt}\end{center}

\subsection{3.
関連の可能性がある先行プログラム（同じ問題意識）}\label{ux95a2ux9023ux306eux53efux80fdux6027ux304cux3042ux308bux5148ux884cux30d7ux30edux30b0ux30e9ux30e0ux540cux3058ux554fux984cux610fux8b58}

ボルン則を、基本公準としてではなく、より深い決定論的層と検出過程から創発／近似として導こうとするプログラムが複数存在する。これらはビート予想と\textbf{問題意識を共有しうる}。一方で、ビート／無理比の等分布／整合比依存の逸脱といったビート予想の要素が、これらと機構として接続するか否かは、現時点では未確定であり、今後の課題である。

\subsubsection{3.1 古典場＋閾値検出からのボルン則創発（La
Cour--Williamson）}\label{ux53e4ux5178ux5834ux95beux5024ux691cux51faux304bux3089ux306eux30dcux30ebux30f3ux5247ux5275ux767ala-courwilliamson}

La Cour と Williamson {[}3{]}（\emph{Quantum} 4, 350,
2020）は、量子光学の文脈で、古典場と「強度閾値型検出器」という決定論的な検出モデル、および実在（virtual
でない）と仮定した QED
真空のみを用いて、量子的とされるいくつかの現象を再現できることを示した。離散的な検出事象に標準的なデータ解析の慣行を適用するだけで量子--古典境界が再現される、という構図である。

問題意識（決定論的基層＋閾値検出からのボルン則の創発）はビート予想と近い可能性がある。ランダム性の起源を観測器--系の振動数ビートに同定するという本予想の論法が本研究と接続するか否かは、本補遺では判断せず、今後の課題とする。

\subsubsection{3.2 プリ量子古典統計場理論（PCSFT,
Khrennikov）}\label{ux30d7ux30eaux91cfux5b50ux53e4ux5178ux7d71ux8a08ux5834ux7406ux8ad6pcsft-khrennikov}

Khrennikov のプリ量子古典統計場理論（PCSFT）{[}4{]}（\emph{Phys. Lett.
A} 372, 6588, 2008；関連して
{[}5{]}）は、量子粒子を「プリ量子ランダム場」の象徴的表現とみなし、古典ランダム場が閾値検出器と相互作用して生じる確率がボルン則に\textbf{近似的にのみ}一致する、と主張する。ボルン則は
PCSFT
から近似理論として現れ、その帰結として\textbf{ボルン則からの逸脱が予言}される。

PCSFT は、ビート予想 {[}1{]} が既に系譜A
として引用している枠組みであり（{[}1{]}
§6.1）、「ボルン則＝決定論的基層への近似・逸脱の予言」という点で本予想と
spirit が近い可能性がある。二点を補足する。第一に、PCSFT
が予言する逸脱が、観測器--系の位相コヒーレンスや整合比に依存する形でビート予想と接続するか否かは、現時点では未確定であり、今後の課題である。第二に、Khrennikov
が {[}4{]}
等で逸脱の傍証として参照した初期の三スリット実験以降、高精度の Sorkin
型検定はいずれもボルン則からの有意な逸脱を検出していない（§4.1
参照）。すなわち PCSFT
の中核予言は現時点で実験的に確認されておらず、本補遺はこれを支持証拠としては扱わない。

\subsubsection{3.3 任意関数の方法（Feintzeig
ら）}\label{ux4efbux610fux95a2ux6570ux306eux65b9ux6cd5feintzeig-ux3089}

Bonds・Burson・Feintzeig ら {[}6{]}（arXiv:2409.16457,
2024）は、ユニタリなシュレディンガー動力学に対し、ある力学パラメータを任意の初期分布をもつ確率変数として扱う「任意関数の方法」を適用し、長時間かつ小
\(\hbar\)
の同時極限でボルン則確率が普遍的極限挙動として現れることを示した。これはビート予想
{[}1{]} が系譜B として引用している枠組みである（{[}1{]} §4, §6.2）。

本予想は、{[}6{]}
が抽象的に置く「等分布する力学パラメータ」を、観測器--系のヘテロダイン相対位相として同定する方向にある点で、{[}6{]}
と関連する可能性がある。{[}6{]}
自体には、実験提案やボルン則からの逸脱予言は、本補遺の確認した範囲では含まれていない。両者が機構として接続するか否かは今後の課題である。

\subsubsection{3.4
共通点と相違点（まとめ）}\label{ux5171ux901aux70b9ux3068ux76f8ux9055ux70b9ux307eux3068ux3081}

§3.1--3.3
の三プログラムは、いずれも「ボルン則を決定論的・古典的な基層と検出過程から創発／近似として導く」という問題意識をビート予想と共有しうる。一方で、\textbf{観測器--系の振動数ビート、無理比に対する
Weyl
等分布、コヒーレント／整合比領域でのボルン則の質的な破れ}という、ビート予想の要素が、これらと機構として接続するか否かは、本補遺では判断せず、今後の課題として残す。複数の独立な道が「ボルン則は基本公準ではないかもしれない」という同じ問いに向かっている可能性があること自体は、記録に値する。

\begin{center}\rule{0.5\linewidth}{0.5pt}\end{center}

\subsection{4.
別方向の研究との関連可能性}\label{ux5225ux65b9ux5411ux306eux7814ux7a76ux3068ux306eux95a2ux9023ux53efux80fdux6027}

\subsubsection{4.1 コヒーレント局部発振器を用いるボルン則検定（Conlon
ら）}\label{ux30b3ux30d2ux30fcux30ecux30f3ux30c8ux5c40ux90e8ux767aux632fux5668ux3092ux7528ux3044ux308bux30dcux30ebux30f3ux5247ux691cux5b9aconlon-ux3089}

Conlon ら {[}7{]}（\emph{New J. Phys.} 26, 053003,
2024）は、コヒーレント光状態を用いた三腕干渉計とホモダイン検出により、ボルン則（三次／Sorkin
干渉）とヒルベルト空間の複素性を検定した。結果は、Sorkin パラメータ
\(\kappa = 0.002 \pm 0.004\)、Peres パラメータ \(F = 1.0000 \pm 0.0003\)
であり、いずれも標準量子力学と整合する（ボルン則からの逸脱は検出されていない）。

本研究は、本補遺が確認した範囲で、コヒーレントな局部発振器（ホモダイン）を実際に用いた数少ないボルン則検定であり、ビート予想と関連する可能性がある。測定対象は高次（Sorkin）干渉であり、ビート予想
{[}1{]}
の判別量――固定した観測器--系相対位相のもとでの単発射影統計が整合比（有理／無理）にどう依存するか――との接続は、本補遺では判断せず今後の課題とする。コヒーレント
LO を用いる Born
検定が既にあること、本補遺の確認した範囲ではそれが本予想の判別領域での測定を含まないことを、ここに記録する。

\subsubsection{4.2 決定論的基層だがボルン則の逸脱に消極的な立場（'t
Hooft）}\label{ux6c7aux5b9aux8ad6ux7684ux57faux5c64ux3060ux304cux30dcux30ebux30f3ux5247ux306eux9038ux8131ux306bux6d88ux6975ux7684ux306aux7acbux5834t-hooft}

't Hooft のセル・オートマトン解釈 {[}8{]}（Springer, \emph{Fundamental
Theories of Physics} 185,
2016）は、量子力学を基礎的な決定論系を扱う道具とみなす。これはビート予想
{[}1{]} が既に引用する隣接系譜である（{[}1{]} §6.3）。

't Hooft
の枠組みは、決定論的基層を置きながらも、再現可能なボルン則の逸脱には消極的とされる（オントロジカルな確率の写像としてボルン則が保たれる方向を志向する）。前論文
v0.4 {[}1{]}
が「より根源的な不確定性により完全コヒーレントな観測が実現されない可能性があり（その機構の存在は仮説であって主張されない）、もしそうなら標準理論物理学のボルン則がより強く成立する」という立場をとることを踏まえると、観測可能なボルン則の逸脱を積極的には主張しないという点で、ビート予想（v0.4）と
't Hooft
は------理由は異なるものの------結論として近い。一方、観測可能な逸脱を予言する
PCSFT（§3.2）は、この点で両者と方向を異にする。同一の「決定論的基層」という問題意識の中にこうした濃淡があることを記録する。

\begin{center}\rule{0.5\linewidth}{0.5pt}\end{center}

\subsection{5.
検証可能領域の所在：精密分光は判別測定を行っていない}\label{ux691cux8a3cux53efux80fdux9818ux57dfux306eux6240ux5728ux7cbeux5bc6ux5206ux5149ux306fux5224ux5225ux6e2cux5b9aux3092ux884cux3063ux3066ux3044ux306aux3044}

ビート予想 {[}1{]}
は、原子時計のラムゼー計量（観測器＝局部発振器と原子の相対位相
\(\varphi_0=(\nu_{\rm atom}-\nu_{\rm LO})\,T\) をヘテロダインで読む）に
operational に接地する。したがって原子時計は、本予想の最も自然な検証
testbed
である。本節は、その文献が判別に必要な測定を行っているかを確認した結果を記録する。

原子時計の標準的な雑音記述では、量子射影雑音（QPN）は、測定の確率的性質（ボルン則的射影）に由来する\textbf{還元不能な二項雑音床}として扱われる
{[}9{]}（Itano ら, \emph{Phys. Rev.~A} 47, 3554, 1993；総説として
{[}10{]} Ludlow ら, \emph{Rev.~Mod. Phys.} 87, 637,
2015）。実機でも、ロック点近傍で QPN
を下回る／非二項的な単発統計が現れるという異常は報告されていない。長時間相互作用やデッドタイム除去によって原子--LO
のコヒーレンスを部分的に高めた構成（{[}11{]} Schioppo ら, \emph{Nature
Photonics} 11, 48, 2017；{[}12{]} Al-Masoudi ら, \emph{Phys. Rev.~A} 92,
063814, 2015）でも、性能は QPN 床に支配され、異常は見られない。

QPN／標準量子限界を下回る実演は、いずれも\textbf{スピンスクイージング（エンタングルメント）によって意図的に作られたもの}であり、エンタングルメント目撃量（Wineland
パラメータ等）で認証されている（例：{[}13{]} Malia ら, \emph{Phys.
Rev.~Lett.} 125, 043202,
2020）。本補遺は、これらをビート予想の傍証としても反例としても扱わない。なお、スクイージングによらない原子時計の制約（Dick
効果、レーザー位相雑音、コヒーレンス時間限界）も古典雑音として記述されている（{[}14{]}
Schulte ら, \emph{Nat. Commun.} 11, 5955, 2020）。

要点は次である。原子時計はオンレゾナンス・位相ロック
LO・射影読み出しという、ビート予想の判別領域のすぐ手前の条件で日常的に運用されているが、判別に必要な測定――固定した原子--LO
相対位相のもとで単発励起率の統計を採り、整合比（有理／無理）の関数として非二項構造を探す測定――は、確認した範囲のいずれの文献でも行われていない。通常の運用は、ディザと雑音により実効的に等分布領域（無理比側）にあり、そこではビート予想と標準量子力学の予言は一致する。両者が分岐する固定コヒーレンス領域は、時計づくりの目的では不要なため、踏み込まれていない。これは隠蔽ではなく、測定プロトコルが構造的にこの領域に入らないことの帰結である。

前論文 v0.4 {[}1{]} は、この \(O(1/N)\)
補正を測定可能な逸脱とは主張せず、より根源的な不確定性が完全コヒーレント（厳密な整数比）な観測の実現そのものを妨げるような機構が存在する可能性があるという仮説を置き（その存在は主張しない）、もしそれがあれば標準理論物理学のボルン則がより強く成立する、という立場をとる。本補遺が確認した「ロック点近傍で異常が報告されていない／判別に必要な測定が行われていない」という事実は、この立場と矛盾しない（逸脱の非観測は前論文の想定と整合的である）。判別測定の有無を確認したことを、本補遺の記録とする。

\begin{center}\rule{0.5\linewidth}{0.5pt}\end{center}

\subsection{6.
まとめ：接続は今後の課題}\label{ux307eux3068ux3081ux63a5ux7d9aux306fux4ecaux5f8cux306eux8ab2ux984c}

本補遺は、ビート予想 {[}1{]} の公開後に行った関連文献調査を記録した。

\begin{itemize}
\tightlist
\item
  ビート予想の中核機構（系--観測器ビート＋無理比 Weyl
  等分布によるボルン統計の創発、コヒーレント／整合比領域での逸脱）を同じ形で提案・検証・観測した先行研究は、調査範囲では確認されなかった。
\item
  一方、「ボルン則を決定論的基層＋閾値検出から創発／近似として導く」という問題意識を共有しうる先行プログラムが存在する（La
  Cour--Williamson {[}3{]}、Khrennikov PCSFT {[}4,5{]}、任意関数の方法
  {[}6{]}）。ビート予想はこれらと関連する可能性があるが、機構として接続するか否かは今後の課題である。
\item
  コヒーレント LO を用いるボルン則検定（Conlon ら
  {[}7{]}）は既に存在する。Sorkin 高次干渉の null
  検定であり、本予想の判別領域での測定との接続は今後の課題である。
\item
  原子時計は本予想の自然な testbed
  だが、判別に必要な測定（固定相対位相での単発射影統計の整合比依存）は、確認した範囲では行われておらず、予言は反証ではなく未検証の状態にある。
\end{itemize}

これらの先行研究とビート予想が論理的に接続するか否か、また判別測定が実際に逸脱を示すか否かは、いずれも\textbf{今後の課題}である。本補遺は、関連の見込まれる範囲で必要部分を引用し、その関連可能性に触れるにとどめる。ビート予想の地位（予想であり、核は条件付きで証明済み、橋は未証明、単発性・指数2・スケール帰属は未解決
{[}1{]}）は、本補遺によって変わらない。

\begin{center}\rule{0.5\linewidth}{0.5pt}\end{center}

\subsection{参考文献}\label{ux53c2ux8003ux6587ux732e}

{[}1{]} 木原 範昭, 「観測者--系ビートの等分布によるボルン統計の創発 ―
局在核モデル上の予想」, Zenodo, v0.4 (2026-06-28), Version DOI:
10.5281/zenodo.20985736 / Concept DOI:
10.5281/zenodo.20967081.（本補遺の対象＝前論文）

{[}2{]} 木原 範昭, 「局在奇数倍音波の再生核性によるボルン分布の形の導出
― 残す公準を二乗則とランダム性に絞る」, Zenodo, Concept DOI:
10.5281/zenodo.20965526.

{[}3{]} B. R. La Cour and M. C. Williamson, ``Emergence of the Born rule
in quantum optics,'' \emph{Quantum} \textbf{4}, 350 (2020). DOI:
10.22331/q-2020-10-26-350.

{[}4{]} A. Khrennikov, ``Detection model based on representation of
quantum particles by classical random fields: Born's rule and beyond,''
\emph{Phys. Lett. A} \textbf{372}, 6588--6592 (2008). DOI:
10.1016/j.physleta.2008.09.027.

{[}5{]} A. Khrennikov, ``Towards violation of Born's rule: description
of a simple experiment,'' \emph{AIP Conf. Proc.} \textbf{1327}, 387--393
(2011). DOI: 10.1063/1.3578726.

{[}6{]} L. Bonds, B. Burson, K. Cicchella, B. H. Feintzeig, Lynnx, A.
Yusaini, ``Quantum Probability via the Method of Arbitrary Functions,''
arXiv:2409.16457 (2024).

{[}7{]} L. O. Conlon, A. Walsh, Y. Hua, O. Thearle, T. Vogl, F.
Eilenberger, P. K. Lam, S. M. Assad, ``Testing the postulates of quantum
mechanics with coherent states of light and homodyne detection,''
\emph{New J. Phys.} \textbf{26}, 053003 (2024). DOI:
10.1088/1367-2630/ad3f1a.

{[}8{]} G. 't Hooft, \emph{The Cellular Automaton Interpretation of
Quantum Mechanics}, Fundamental Theories of Physics \textbf{185}
(Springer, 2016). DOI: 10.1007/978-3-319-41285-6.

{[}9{]} W. M. Itano, J. C. Bergquist, J. J. Bollinger, J. M. Gilligan,
D. J. Heinzen, F. L. Moore, M. G. Raizen, D. J. Wineland, ``Quantum
projection noise: Population fluctuations in two-level systems,''
\emph{Phys. Rev.~A} \textbf{47}, 3554 (1993). DOI:
10.1103/PhysRevA.47.3554.

{[}10{]} A. D. Ludlow, M. M. Boyd, J. Ye, E. Peik, P. O. Schmidt,
``Optical atomic clocks,'' \emph{Rev.~Mod. Phys.} \textbf{87}, 637
(2015). DOI: 10.1103/RevModPhys.87.637.

{[}11{]} M. Schioppo, R. C. Brown, W. F. McGrew, N. Hinkley, et al.,
``Ultrastable optical clock with two cold-atom ensembles,'' \emph{Nature
Photonics} \textbf{11}, 48--52 (2017). DOI: 10.1038/nphoton.2016.231.

{[}12{]} A. Al-Masoudi, S. Dörscher, S. Häfner, U. Sterr, C. Lisdat,
``Noise and instability of an optical lattice clock,'' \emph{Phys.
Rev.~A} \textbf{92}, 063814 (2015). DOI: 10.1103/PhysRevA.92.063814.

{[}13{]} B. K. Malia, J. Martínez-Rincón, Y. Wu, O. Hosten, M. A.
Kasevich, ``Free space Ramsey spectroscopy in rubidium with noise below
the quantum projection limit,'' \emph{Phys. Rev.~Lett.} \textbf{125},
043202 (2020). DOI: 10.1103/PhysRevLett.125.043202.

{[}14{]} M. Schulte, C. Lisdat, P. O. Schmidt, U. Sterr, K. Hammerer,
``Prospects and challenges for squeezing-enhanced optical atomic
clocks,'' \emph{Nat. Commun.} \textbf{11}, 5955 (2020). DOI:
10.1038/s41467-020-19403-7.

{[}15{]} H. Weyl, ``Über die Gleichverteilung von Zahlen mod. Eins,''
\emph{Math. Ann.} \textbf{77}, 313--352 (1916). DOI: 10.1007/BF01475864.

\begin{center}\rule{0.5\linewidth}{0.5pt}\end{center}

\emph{(本補遺は、前論文 {[}1{]}
公開後の関連文献調査の記録である。本補遺は機構の同一性も実験的確認も主張せず、関連の見込まれる範囲でその可能性に触れるにとどめる。前論文との論理的接続、および判別測定の実施は、今後の課題として残す。)}

\end{document}
